\section{Design of SAMM}
\label{sec:design}
\label{sec:data_organization}
\label{sec:samm}

\subsection{Architecture Overview}
\label{sec:samm-arch}

SAMM is a semantic overlay allocator. It interposes a slot-aware region-management layer between the application and a commodity allocator (e.g., jemalloc~\cite{jemalloc}), without replacing it. The architecture is hierarchical, as shown in Figure~\ref{fig:samm-arch}.

\newsavebox{\sammArchBox}\sbox{\sammArchBox}{% Fig. -- SAMM architecture: the layers between the KVS and the operating system.
% Same conventions as figs/1.tex: 252-unit-wide grid (x unit = \columnwidth/252),
% y in pt growing downwards, labels at 7.5pt Helvetica. Grayscale only.
\begin{tikzpicture}[x={\columnwidth/252},y=-1pt,
    font=\fontsize{7.5}{9}\fontfamily{phv}\selectfont,
    every node/.append style={text height=5.4pt,text depth=1.6pt},
    >={Stealth[length=4pt,width=3.2pt]},
    ours/.style={draw=ink,line width=0.6pt,fill=ink!30,rounded corners=1.5pt},
    base/.style={draw=ink,line width=0.6pt,fill=ink!8,rounded corners=1.5pt},
    app/.style={draw=ink!70,line width=0.6pt,fill=white,rounded corners=1.5pt,dash pattern=on 2pt off 1.4pt},
    call/.style={->,draw=ink,line width=0.6pt},
    lbl/.style={inner sep=1pt,text=ink!80,font=\fontsize{7}{8}\fontfamily{phv}\itshape},
  ]
  \definecolor{ink}{RGB}{34,34,34}

  % layers (top to bottom)
  \filldraw[app]  (63,4)   rectangle (189,20);  \node at (126,12)  {KVS (hash index)};
  \filldraw[ours] (63,34)  rectangle (189,50);  \node at (126,42)  {segment allocator};
  \filldraw[ours] (63,50)  rectangle (189,66);  \node at (126,58)  {slot table and block lists};
  \filldraw[base] (63,80)  rectangle (189,96);  \node at (126,88)  {jemalloc};
  \filldraw[base] (63,110) rectangle (189,126); \node at (126,118) {OS: mmap, hugepages};

  % the two SAMM layers, grouped by a brace on the left
  \draw[draw=ink!70,line width=0.5pt,decorate,decoration={brace,amplitude=3pt}] (58,65.5) -- (58,34.5);
  \node[inner sep=0pt,font=\fontsize{7.5}{9}\fontfamily{phv}\bfseries] at (38,45.5) {SAMM};
  \node[lbl,inner sep=0pt] at (38,54.5) {(Fig.~\ref{fig:samm})};

  % calls between layers (labels to the right of the arrows)
  \draw[call] (126,20) -- (126,34);  \node[lbl,anchor=west] at (130,27)  {alloc(slot, size)};
  \draw[call] (126,66) -- (126,80);  \node[lbl,anchor=west] at (130,73)  {alloc / free a block};
  \draw[call] (126,96) -- (126,110); \node[lbl,anchor=west] at (130,103) {map memory};

\end{tikzpicture}%
}
\newsavebox{\sammLayoutBox}\sbox{\sammLayoutBox}{\resizebox{0.8\columnwidth}{!}{% Fig. -- SAMM: slot table, per-slot block lists, the layout of one block,
% and the layout of one segment. Grayscale only.
% Same conventions as figs/1.tex: 252-unit-wide grid (x unit = \columnwidth/252),
% y in pt growing downwards, labels at 7.5pt Helvetica.
\begin{tikzpicture}[x={\columnwidth/252},y=-1pt,
    font=\fontsize{7.5}{9}\fontfamily{phv}\selectfont,
    every node/.append style={text height=5.4pt,text depth=1.6pt},
    >={Stealth[length=4pt,width=3.2pt]},
    cellrow/.style={draw=ink,line width=0.6pt,fill=white},
    cellgap/.style={draw=ink,line width=0.6pt,fill=ink!8},
    blk/.style={draw=ink,line width=0.6pt,fill=ink!16,rounded corners=1.5pt,
                minimum width=36pt,minimum height=12pt,inner sep=0pt},
    blkz/.style={blk,draw=ink,line width=1pt,fill=ink!30},
    meta/.style={draw=ink,line width=0.6pt,fill=ink},
    used/.style={draw=ink,line width=0.6pt,fill=ink!16},
    hdr/.style={draw=ink,line width=0.6pt,fill=ink!30},
    free/.style={draw=ink!55,line width=0.6pt,fill=white,dash pattern=on 1.6pt off 1.1pt},
    ptr/.style={->,draw=ink,line width=0.6pt},
    dlink/.style={->,draw=ink!60,line width=0.6pt},
    zoom/.style={fill=ink!6,draw=none},
    ttl/.style={anchor=base west,inner sep=0pt,font=\fontsize{7.5}{9}\fontfamily{phv}\bfseries},
    note/.style={inner sep=1pt,text=ink!80,font=\fontsize{7}{8}\fontfamily{phv}\itshape},
    brc/.style={draw=ink!70,line width=0.5pt,decorate,decoration={brace,amplitude=3pt,mirror}},
  ]
  \definecolor{ink}{RGB}{34,34,34}
  \def\kv#1{\textit{k}\textsubscript{#1}, \textit{v}\textsubscript{#1}}

  % titles
  \node[ttl,anchor=base] at (26,6) {Slot table};
  \node[ttl] at (70,6) {Blocks of each slot};

  % zoom areas, drawn first so that everything else sits on top
  \path[zoom] (126.5,53) -- (163.5,53) -- (248,92) -- (4,92) -- cycle;
  \path[zoom] (94,110) -- (150,110) -- (212,128) -- (40,128) -- cycle;

  % slot table
  \filldraw[cellrow] (4,12) rectangle (48,26);
  \filldraw[cellgap] (4,26) rectangle (48,40);
  \filldraw[cellrow] (4,40) rectangle (48,54);
  \filldraw[cellgap] (4,54) rectangle (48,68);
  \node at (26,19) {slot \textit{i}};
  \node[text=ink!50] at (26,33) {$\cdots$};
  \node at (26,47) {slot \textit{j}};
  \node[text=ink!50] at (26,61) {$\cdots$};

  % blocks of slot i
  \draw[ptr] (48,19) -- (70,19);
  \node[blk] (bi0) at (90,19)  {block 0};
  \node[blk] (bi1) at (145,19) {block 1};
  \node[blk] (bi2) at (200,19) {block 2};
  \draw[dlink] ([yshift=2.2pt]bi0.east) -- ([yshift=2.2pt]bi1.west);
  \draw[dlink] ([yshift=-2.2pt]bi1.west) -- ([yshift=-2.2pt]bi0.east);
  \draw[dlink] ([yshift=2.2pt]bi1.east) -- ([yshift=2.2pt]bi2.west);
  \draw[dlink] ([yshift=-2.2pt]bi2.west) -- ([yshift=-2.2pt]bi1.east);

  % blocks of slot j; block 1 is shown enlarged below
  \draw[ptr] (48,47) -- (70,47);
  \node[blk]  (bj0) at (90,47)  {block 0};
  \node[blkz] (bj1) at (145,47) {block 1};
  \draw[dlink] ([yshift=2.2pt]bj0.east) -- ([yshift=2.2pt]bj1.west);
  \draw[dlink] ([yshift=-2.2pt]bj1.west) -- ([yshift=-2.2pt]bj0.east);

  % one block: prologue, segments, epilogue
  \filldraw[meta] (4,92) rectangle (22,110);
  \node[text=white] at (13,101) {P};
  \filldraw[used] (22,92) rectangle (70,110);
  \filldraw[free] (70,92) rectangle (94,110);
  \filldraw[used] (94,92) rectangle (150,110);
  \filldraw[free] (150,92) rectangle (174,110);
  \filldraw[used] (174,92) rectangle (230,110);
  \filldraw[meta] (230,92) rectangle (248,110);
  \node[text=white] at (239,101) {E};
  \node at (46,101)  {\kv{1}};
  \node at (122,101) {\kv{2}};
  \node at (202,101) {\kv{3}};
  \node[text=ink!55] at (82,101)  {free};
  \node[text=ink!55] at (162,101) {free};

  % labels under the block
  \node[note] at (13,118) {prologue};
  \node[note] at (239,118) {epilogue};

  % one segment: header fields, then the payload
  \filldraw[hdr]  (40,128) rectangle (68,142);
  \filldraw[hdr]  (68,128) rectangle (96,142);
  \filldraw[used] (96,128) rectangle (140,142);
  \filldraw[used] (140,128) rectangle (212,142);
  \node at (54,135)  {size};
  \node at (82,135)  {flags};
  \node at (118,135) {key};
  \node at (176,135) {value};
  \draw[brc] (40.5,146) -- (94,146);
  \draw[brc] (98,146) -- (211.5,146);
  \node[note] at (67,156)  {header};
  \node[note] at (155,156) {payload};

  % pointer that SAMM returns to the KVS
  \draw[ptr] (96,165) -- (96,145);
  \node[note,anchor=north] at (96,166) {pointer returned to the KVS};
  % tight bounding box (no empty band under the last label)
  \pgfresetboundingbox
  \path (0,-2) rectangle (252,175);
\end{tikzpicture}%
}}
\newlength{\sammFigH}\setlength{\sammFigH}{\dimexpr\ht\sammLayoutBox+\dp\sammLayoutBox\relax}

\begin{figure}[t]
  \centering
  \parbox[c][\sammFigH][c]{\columnwidth}{\centering\usebox{\sammArchBox}}
  \caption{SAMM four-layer architecture. Jemalloc provides coarse-grained blocks. SAMM enforces slot ownership and manages intra-block segments.}
  \label{fig:samm-arch}
\end{figure}

At the lowest layer, jemalloc supplies large aligned memory blocks. SAMM acquires these blocks, tags each one with a slot owner, and registers it in a central slot table. No block is ever shared between two slots. Within each block, SAMM manages variable-sized segments, one per key-value pair. Allocation, deallocation, and segment reuse are governed by a per-slot free list that threads across all blocks belonging to that slot.

\subsection{Memory Layout}
\label{sec:samm-layout}
\label{sec:samm-blocks}
\label{sec:samm-block-layout}
\label{sec:samm-segments}

\begin{figure}[t]
  \centering
  \parbox[c][\sammFigH][c]{\columnwidth}{\centering\usebox{\sammLayoutBox}}
  \caption{SAMM memory layout. Each slot owns a doubly linked list of blocks, and each block holds variable-size segments of that slot only, between a prologue and an epilogue. Free segments of a slot form one free list that can span its blocks.}
  \label{fig:samm}
\end{figure}

\paragraph{Slot table.}
The slot table is a flat array indexed by slot identifier. Each entry represents one keyspace partition and tracks:

The slot table entry is the sole authoritative index of a slot's memory footprint. During slot migration, SAMM traverses the block list starting at head and emits each block's body as a contiguous byte range. No scatter-gather over the global heap is required.

\paragraph{Block layout.}
A block is a fixed-size, contiguous memory block obtained from jemalloc (e.g., 2\,MB, aligned to a hugepage boundary). Each block is exclusively owned by one slot. Rather than a slotted-page directory, SAMM structures each block with a prologue and an epilogue---fixed-size metadata blocks at the low and high ends of the block respectively. The body between them holds a dense, append-style sequence of variable-sized segments (Fig.~\ref{fig:samm}).

\paragraph{Segment layout.}
A segment is the atomic unit of storage for a single key-value pair. Segments are variable-sized. Each segment is exactly large enough to hold the fixed-size header plus the key and value payloads, with no wasted padding beyond alignment. The segment format is:

The free list is implemented as an intrusive linked list that overlays the first 8 bytes of the payload region. When \texttt{FLAG\_FREE} is set, those bytes are reinterpreted as a \texttt{next\_free} pointer chaining this segment to the next free segment in the slot's free list, which may reside in any of the slot's blocks without any auxiliary data structure.

On allocation, SAMM returns to the caller a direct pointer to the first byte of the payload region immediately following the segment header. The caller writes the key and value directly into the block at that address, incurring no intermediate copy.

When SAMM needs a new block for a slot, it invokes the arena's alloc hook, receiving a hugepage-aligned region that it then initialises with a prologue and epilogue before linking it to the slot's block list. When a slot's block becomes empty (all segments freed), SAMM unlinks it from the doubly-linked list and returns it to jemalloc via the dalloc hook. This delegation avoids reimplementing OS-level concerns such as NUMA-aware placement, transparent hugepage promotion, and background purging.

\section{Fault-tolerant migration protocol}
\label{sec:protocol}

\begin{table}[t]
  \centering
  \footnotesize
  \renewcommand{\arraystretch}{1.2}
  \setlength{\tabcolsep}{4pt}
  \begin{tabular*}{\columnwidth}{@{}l@{\extracolsep{\fill}}cl@{}}
    \toprule
    \multicolumn{1}{c}{\textbf{Entry}} & \textbf{Log} & \multicolumn{1}{c}{\textbf{Use}} \\
    \midrule
    \texttt{START($G_i$)}         & donor     & resume $G_i$ after a crash \\
    \texttt{TXN\_START($G_i$)}    & recipient & keep tombstones for $G_i$ \\
    \texttt{TRANSFER\_DURABLE($G_i$)} & recipient & stop double reads for $G_i$ \\
    \texttt{DONE($G_i$)}          & donor     & skip $G_i$ on recovery \\
    \bottomrule
  \end{tabular*}
  \caption{Replicated-log entries of group $G_i$.}
  \label{tab:log-entries}
\end{table}

\begin{figure}[t]
  \centering
  \begin{tikzpicture}[x={\columnwidth/252},y=-1pt,
    font=\fontsize{7.5}{9}\fontfamily{phv}\selectfont,
    >={Stealth[length=4pt,width=3.2pt]},
    lifeline/.style={draw=lifegray,line width=0.8pt},
    shard/.style={draw=ink,line width=0.6pt,fill=white,rounded corners=1.5pt,
                  minimum width=24pt,minimum height=12pt,inner sep=0pt},
    leader/.style={shard,fill=ink,text=white,font=\fontsize{7.5}{9}\fontfamily{phv}\bfseries},
    msg/.style={->,draw=ink,line width=0.6pt},
    bulk/.style={->,draw=ink,line width=1.8pt,>={Stealth[length=5pt,width=5pt]}},
    notify/.style={->,draw=ink,line width=0.6pt,dash pattern=on 2.2pt off 1.4pt},
    lbl/.style={inner sep=1pt,text=ink},
    phase/.style={rounded corners=2pt},
    phaselbl/.style={anchor=north west,inner sep=1.5pt,outer sep=0pt,
                     font=\fontsize{7.5}{9}\fontfamily{phv}\bfseries},
  ]
  \definecolor{ink}{RGB}{34,34,34}
  \definecolor{lifegray}{RGB}{165,165,165}
  \definecolor{llfill}{gray}{0.93}
  \definecolor{llink}{gray}{0.13}
  \definecolor{irfill}{gray}{0.86}
  \definecolor{irink}{gray}{0.13}
  \def\crashat#1#2#3#4{%
    \draw[ink,line width=1.2pt,line cap=round] (#1-3.2,#2-3.2) -- (#1+3.2,#2+3.2) (#1-3.2,#2+3.2) -- (#1+3.2,#2-3.2);
    \ifx#3e\node[anchor=east,inner sep=0pt,font=\fontsize{6.6}{7}\fontfamily{phv}\bfseries] at (#1-5,#2) {#4};
    \else\node[anchor=west,inner sep=0pt,font=\fontsize{6.6}{7}\fontfamily{phv}\bfseries] at (#1+5,#2) {#4};\fi}
  \def\Gi{\textit{G}\textsubscript{\textit{i}}}
  \def\us{{\fontencoding{T1}\selectfont\_}}%
  \def\fstep#1{\tikz[baseline={([yshift=-2.7pt]s.center)}]\node[circle,fill=ink,text=white,inner sep=0pt,minimum size=9.6pt,text height=5.1pt,text depth=0.7pt,font=\fontsize{7.2}{7.2}\fontfamily{phv}\bfseries\selectfont](s){\ifnum#1=1 \kern0.84pt\fi#1};}

  \def\Ao{22}  \def\At{62}  \def\Ah{104}
  \def\Bo{150} \def\Bt{193} \def\Bh{232}
  \def\bottom{340}

  \fill[phase,llfill] (0,61) rectangle (252,146);
  \fill[phase,irfill] (0,151) rectangle (252,265);
  \fill[phase,irfill] (0,285) rectangle (252,333);

  \foreach \x in {\Ao,\At,\Ah,\Bo,\Bt,\Bh} \draw[lifeline] (\x,33) -- (\x,\bottom);

  \node[font=\fontsize{7.5}{9}\fontfamily{phv}\bfseries,anchor=base] at (\At,6) {Donor shard group};
  \node[font=\fontsize{7.5}{9}\fontfamily{phv}\bfseries,anchor=base] at (\Bt,6) {Recipient shard group};
  \draw[ink,line width=0.5pt] (\Ao-12,15) -- ++(0,-3) -- (\Ah+14,12) -- ++(0,3);
  \draw[ink,line width=0.5pt] (\Bo-14,15) -- ++(0,-3) -- (\Bh+12,12) -- ++(0,3);
  \node[shard]  at (\Ao,27) {A1};
  \node[shard]  at (\At,27) {A2};
  \node[leader] at (\Ah,27) {A3};
  \node[leader] at (\Bo,27) {B1};
  \node[shard]  at (\Bt,27) {B2};
  \node[shard]  at (\Bh,27) {B3};
  \node[lbl,fill=white,font=\fontsize{7}{8}\fontfamily{phv}\itshape] at (\Ah,41) {leader};
  \node[lbl,fill=white,font=\fontsize{7}{8}\fontfamily{phv}\itshape] at (\Bo,41) {leader};

  \node[lbl,above=2pt,fill=white] at ({(\Ao+\Ah)/2},54) {\fstep{1}\hspace{2pt}START(\Gi)};
  \draw[notify] (\Ah,54) -- (\Ao,54);
  \draw[notify] (\Ah,54) -- (\At,54);

  \node[phaselbl,text=llink,fill=llfill] at (2,63) {Leader-to-leader};
  \draw[msg]  (\Ah,82) -- node[lbl,above=2pt] {\fstep{2}\hspace{2pt}prepare}  (\Bo,82);
  \node[lbl,above=2pt,fill=llfill] at ({(\Bo+\Bh)/2},95) {TXN\us START(\Gi)};
  \draw[notify] (\Bo,95) -- (\Bt,95);
  \draw[notify] (\Bo,95) -- (\Bh,95);
  \draw[msg]  (\Bo,108) -- node[lbl,above=2pt] {MRs}      (\Ah,108);
  \draw[bulk] (\Ah,123) -- node[lbl,above=2pt] {\fstep{3}\hspace{2pt}transfer} (\Bo,123);
  \draw[msg]  (\Ah,136) -- node[lbl,above=2pt,fill=llfill,font=\fontsize{6.6}{7}\fontfamily{phv}\selectfont,inner sep=0.6pt] {RDMA\us DONE}     (\Bo,136);
  \draw[msg]  (\Bo,141) -- (\Ah,141);

  \node[phaselbl,text=irink,fill=irfill] at (2,153) {Intra-recipient};
  \draw[msg,rounded corners=2pt] (\Bo,162.5) -- ++(12,0) -- ++(0,9) -- (\Bo,171.5);
  \node[lbl,right=3pt,fill=irfill] at (\Bo+12,167) {\fstep{4}\hspace{2pt}ack B1 (1/3)};
  \draw[bulk] (\Bo,193) -- node[lbl,above=2pt,align=center,font=\fontsize{7.5}{8}\fontfamily{phv}\selectfont] {chain\\replication} (\Bt,193);
  \draw[msg]  (\Bo,207) -- node[lbl,above=2pt,fill=irfill,font=\fontsize{6.6}{7}\fontfamily{phv}\selectfont,inner sep=0.6pt] {CHAIN\us DONE}     (\Bt,207);
  \draw[msg]  (\Bt,212) -- (\Bo,212);
  \draw[msg]  (\Bt,227) -- node[lbl,above=2pt,fill=irfill] {ack B2 (2/3)} (\Bo,227);
  \draw[ink,line width=0.6pt,decorate,decoration={brace,amplitude=4pt}]
    (\Bo-7,231) -- node[lbl,left=5pt,align=center] {IU\\(\S\ref{sec:index_updating})} (\Bo-7,158);

  \draw[irink,line width=1pt] (\Bo,241) -- (\Bh,241)
    node[midway,lbl,inner sep=2pt,fill=irfill,text=irink,
         font=\fontsize{7.5}{9}\fontfamily{phv}\bfseries] {majority reached};
  \node[inner sep=0pt,anchor=east] at (\Bo-3,241) {\fstep{5}};
  \node[lbl,above=2pt,anchor=south west,fill=irfill] at (\Bo+3,259) {TRANSFER\us DURABLE(\Gi)};
  \node[inner sep=0pt,anchor=east] at (\Bo-3,253) {\fstep{6}};
  \draw[notify] (\Bo,259) -- (\Bt,259);
  \draw[notify] (\Bo,259) -- (\Bh,259);

  \node[lbl,above=2pt,fill=white] at ({(\Ao+\Ah)/2},279) {\fstep{7}\hspace{2pt}DONE(\Gi)};
  \draw[msg] (\Bo,279) -- (\Ah,279);
  \draw[notify] (\Ah,279) -- (\Ao,279);
  \draw[notify] (\Ah,279) -- (\At,279);

  \node[phaselbl,text=irink,fill=irfill] at (2,287) {Intra-recipient (slow path)};
  \draw[bulk] (\Bt,308) -- node[lbl,above=2pt,align=center,font=\fontsize{7.5}{8}\fontfamily{phv}\selectfont] {chain\\replication}  (\Bh,308);
  \draw[msg]  (\Bt,322) -- node[lbl,above=2pt,fill=irfill,font=\fontsize{6.6}{7}\fontfamily{phv}\selectfont,inner sep=0.6pt] {CHAIN\us DONE}     (\Bh,322);
  \draw[msg]  (\Bh,327) -- (\Bt,327);

  \def\key#1{\tikz[baseline=-2.7pt]\draw[#1] (0pt,0pt) -- (18pt,0pt);\hspace{4pt}}
  \node[inner sep=0pt,anchor=base] at (126,356) {%
    \key{msg}RPC\hspace{16pt}%
    \key{notify}Replicated log entry\hspace{16pt}%
    \key{bulk}RDMA write};
\end{tikzpicture}%

  \caption{Fast path of the migration of one group $G_i$.}%
  \label{fig:migration_flow}
\end{figure}

\subsection{Failure model}
\label{sec:applicability}
\label{sec:fault-model}

\OurImplementation{}'s fault-tolerant migration protocol targets in-memory key-value stores.
We consider a KVS that (i) partitions the key space into logical units such as shards or slots, (ii) replicates each partition across a set of replicas using a leader-based replication protocol~\cite{raft,paxos-simple,chain-replication} with synchronous writes, and (iii) runs on replicas that fail by crashing (\textcolor{blue}{\emph{fail-stop}}).

We call the replicas of a partition a \emph{shard group}. One replica acts as the leader and the others as followers. The leader appends writes and other commands to the group's \emph{replicated log}, and an entry is committed only after it has been durably stored by a majority of the shard group.
A crashed replica loses its in-memory state. It may later rejoin its shard group as a new member and obtain its state from the other replicas through the replication protocol.
Failed replicas are eventually detected by the others, e.g., through timeouts or broken connections.
At most $f$ of the $2f+1$ replicas of each shard group fail. Every committed update, and the data that is being migrated, therefore survives on at least one replica.

An elasticity operation involves two shard groups, a \emph{donor} shard group that currently owns the data being migrated, and a \emph{recipient} shard group that will assume ownership.
Migration proceeds concurrently with client requests.

\subsection{Overview}
\label{sec:migration-protocol}
\label{sec:log-entries}

A migration moves a set of slots from a \emph{donor} shard group to a \emph{recipient} shard group.
A coordinator $C$ splits the slots into \emph{migration groups} $G_1, \dots, G_n$ and instructs the donor leader to migrate them to the recipient leader. Together, the groups form one migration transaction, driven on each leader by the \emph{Migration Manager (MM)}, a dedicated thread that manages the migration.
Each group $G_i = \{s_1, \dots, s_m\}$ is a set of slots, and each slot $s_j$ contains the keys that map to it (e.g., by hashing), together with their values.
In memory, the key-value pairs of a slot are stored in the SAMM blocks that the slot owns (Section~\ref{sec:samm}), so migrating $G_i$ means transferring the blocks of its slots.
Bulk data moves with one-sided RDMA writes and never goes through the log. Only the decisions that must survive a crash are appended to each shard group's replicated log.
Under the fault model of Section~\ref{sec:fault-model}, our protocol ensures that
(i)~committed writes are never lost,
(ii)~a read never returns data that is only partially migrated or not yet replicated, and
(iii)~once a migration starts, it always runs to completion and is never rolled back.
Properties (i) and (ii) come from Fulva~\cite{fulva}.
The protocol has a \emph{fast path}, which a migration follows when no node crashes (Section~\ref{sec:migration-phases}), and recovery rules that resume it after a crash (Section~\ref{sec:failure-handling}).

The progress of each group $G_i$ is recorded by four entries in the donor's and the recipient's replicated logs (Table~\ref{tab:log-entries}).
Every replica learns how far the migration got by applying its log, so a newly elected leader resumes from the last committed entry instead of relying on in-memory state, which a crash loses.
The protocol rests on one invariant:
\emph{if \texttt{TRANSFER\_DURABLE($G_i$)} is committed, a majority of the recipient shard group holds the data of $G_i$.}
It holds because the recipient leader proposes this entry only after a majority of replicas have reported receiving the data (step~\mstep{5}).
Since the fault model keeps a majority of the shard group available and any two majorities intersect, a newly elected leader that lacks part of the data can fetch it from its shard group before serving $G_i$.

\iffalse
\begin{figure*}[t]
  \centering
  % Slot lifecycle -- the states of the slots of one group G_i, left to right,
% and what causes each transition.
% Two-column figure: 504-unit-wide grid (x unit = \textwidth/504), y in pt
% growing downwards, labels at 7.5pt as in figs/1.tex. \input it in a figure*.
\begin{tikzpicture}[x={\textwidth/504},y=-1pt,
    font=\fontsize{7.5}{9}\fontfamily{phv}\selectfont,
    >={Stealth[length=4pt,width=3.2pt]},
    state/.style={draw=ink,line width=0.6pt,fill=white,rounded corners=1.5pt,
                  minimum width=50pt,minimum height=22pt,inner sep=0pt,align=center,
                  font=\fontsize{7.5}{8.5}\fontfamily{phv}\bfseries},
    logged/.style={->,draw=ink,line width=0.6pt,dash pattern=on 2.2pt off 1.4pt},
    soft/.style={->,draw=ink,line width=0.6pt},
    cause/.style={anchor=north,inner sep=0pt,text=ink,align=center,
                  font=\fontsize{7.5}{8.5}\fontfamily{phv}\selectfont},
    serves/.style={anchor=base,inner sep=0pt,
                   font=\fontsize{7.5}{9}\fontfamily{phv}\bfseries},
  ]
  \definecolor{ink}{RGB}{34,34,34}
  \def\Gi{\textit{G}\textsubscript{\textit{i}}}
  \def\us{{\fontencoding{T1}\selectfont\_}}% thin underscore glyph
  \def\eff#1{{\fontshape{sl}\selectfont #1}}% effect of a transition

  % seven columns, 72 units apart; state boxes are 50pt wide
  \def\sy{31}   % y of the state row
  \def\cy{47}   % top of the transition labels

  % who serves G_i in each state
  \node[serves] at (36,6)  {Donor serves \Gi};
  \node[serves] at (180,6) {Writes at recipient, double reads};
  \node[serves] at (396,6) {Recipient serves \Gi{} alone};
  \draw[ink,line width=0.5pt] (11,15)  -- ++(0,-3) -- (61,12)  -- ++(0,3);
  \draw[ink,line width=0.5pt] (83,15)  -- ++(0,-3) -- (277,12) -- ++(0,3);
  \draw[ink,line width=0.5pt] (299,15) -- ++(0,-3) -- (493,12) -- ++(0,3);

  % states
  \node[state] at (36,\sy)  {Owned by\\donor};
  \node[state] at (108,\sy) {Flipped};
  \node[state] at (180,\sy) {Landed};
  \node[state] at (252,\sy) {Chain-\\durable};
  \node[state] at (324,\sy) {Committed};
  \node[state] at (396,\sy) {Merged on\\leader};
  \node[state] at (468,\sy) {Merged on\\followers};

  % transitions
  \draw[logged] (62,\sy)  -- (82,\sy);
  \draw[soft]   (134,\sy) -- (154,\sy);
  \draw[soft]   (206,\sy) -- (226,\sy);
  \draw[logged] (278,\sy) -- (298,\sy);
  \draw[logged] (350,\sy) -- (370,\sy);
  \draw[logged] (422,\sy) -- (442,\sy);

  % what causes each transition, and its effect, under its arrow
  \node[cause] at (72,\cy)  {START(\Gi), then\\TXN\us START(\Gi)\\\eff{new writes go}\\\eff{to the recipient}};
  \node[cause] at (144,\cy) {RDMA\us DONE\\from the donor\\\eff{data is in leader's}\\\eff{landing buffers}};
  \node[cause] at (216,\cy) {CHAIN\us DONE,\\$f$ followers report\\\eff{a majority}\\\eff{holds the data}};
  \node[cause] at (288,\cy) {INDEX\us UPDATE(\Gi)\\commits\\\eff{clients stop}\\\eff{double reads}};
  \node[cause] at (360,\cy) {leader applies\\INDEX\us UPDATE(\Gi)\\\eff{its index points to}\\\eff{the migrated data}};
  \node[cause] at (432,\cy) {followers apply\\INDEX\us UPDATE(\Gi)\\\eff{any of them can}\\\eff{serve \Gi{} as leader}};

  % legend: one centred row
  \def\key#1{\tikz[baseline=-2.7pt]\draw[#1] (0pt,0pt) -- (18pt,0pt);\hspace{4pt}}
  \node[inner sep=0pt,anchor=base] at (252,98) {%
    \key{logged}Replicated log entry\hspace{16pt}%
    \key{soft}Soft state (not logged)};
\end{tikzpicture}%

  \caption{Lifecycle of the slots of one group $G_i$.
  Dashed transitions are caused by a replicated-log entry and survive a crash. Solid ones are soft state.
  \texttt{START($G_i$)} and \texttt{DONE($G_i$)} are in the donor's log, the other entries in the recipient's.
  \texttt{DONE($G_i$)} follows the commit and only closes the group for recovery.}
  \label{fig:slot_lifecycle}
\end{figure*}
\fi

\subsection{Migrating replication protocol}
\label{sec:migration-phases}
\label{sec:intra-donor}
\label{sec:leader-to-leader}
\label{sec:intra-recipient}
\label{sec:commit-completion}

Fig.~\ref{fig:migration_flow} shows the fast path of one group.
The transaction completes once every group has completed.

The donor leader appends \texttt{START($G_i$)} to its log and commits it, making the migration intent durable~(step~\mstep{1}).
It then collects the blocks of the slots of $G_i$ from its slot table (Section~\ref{sec:samm}) and registers them with its RNIC as the source of the transfer, unless they are already registered (e.g., at block allocation, or as part of one memory region that covers the whole address space).

The \emph{leader-to-leader} phase follows.
The donor leader sends a \texttt{prepare} RPC to the recipient leader, which allocates landing buffers for $G_i$, registers them as memory regions, appends \texttt{TXN\_START($G_i$)} to its log, and, once it commits, replies with these memory regions, namely their addresses, sizes, and RDMA access keys (RKeys)~(step~\mstep{2}).
The allocation and registration can also happen out of the critical path, before the \texttt{prepare} arrives, so that the recipient only reports memory regions that are already registered.
From this point, writes to $G_i$ are directed to the recipient shard group.
The donor leader lets the writes to $G_i$ that are already in progress finish and applies no further writes to $G_i$.
Every write that reaches the recipient therefore finds \texttt{TXN\_START($G_i$)} committed, so recipient replicas keep tombstones for it, and the donor's copy of $G_i$ no longer changes during the transfer.
Using these keys, the donor leader writes the data of $G_i$ directly into the landing buffers with one-sided RDMA writes, which do not involve the recipient's CPU, and then sends an \texttt{RDMA\_DONE} RPC to the recipient leader~(step~\mstep{3}).

The \emph{intra-recipient} phase starts on \texttt{RDMA\_DONE}.
Once a majority of replicas, counting the leader, hold the data and the leader's index update has finished~(step~\mstep{5}),
the leader appends \texttt{TRANSFER\_DURABLE($G_i$)} to its replicated log~(step~\mstep{6}).

Finally, the recipient leader notifies the donor leader that $G_i$ is complete, repeating the notification until the donor acknowledges it, and the donor leader appends and commits \texttt{DONE($G_i$)}~(step~\mstep{7}).
Ownership of the slots moves per group, not once per transaction. When \texttt{TRANSFER\_DURABLE($G_i$)} commits, clients stop double reads for $G_i$ and the recipient shard group alone serves $G_i$, while later groups of the same transaction may still be migrating.

\subsection{Index update}
\label{sec:index_updating}

The KVS keeps a dictionary-like index with pointers to the key-value pairs, which must be updated to point to the new memory locations.
This \emph{index update} runs on dedicated background threads, to minimize interference with client-serving threads, and uses fine-grained locking for dictionary accesses, since these threads run concurrently with client threads.

Because new writes for $G_i$ go to the recipient during the migration, the transferred data may contain older versions of these keys. The index update detects and discards such stale entries, keeping only the most recent value of each key.
A key that a client deletes at the recipient during the migration could come back when the index update inserts the donor's older copy of it.
To prevent this, recipient replicas keep a tombstone for every key that is deleted while \texttt{TXN\_START($G_i$)} is active, and the index update skips keys that have a tombstone. The tombstones of $G_i$ are removed once its index update finishes.

\subsection{Serving client requests during migration}
\label{sec:serving-during-migration}

\OurImplementation{} migrates data without interrupting client-facing operations. Requests for slots that are not under migration are handled normally by their owning shard group.
As in the underlying replication protocol, all client requests for $G_i$ are served by shard-group leaders. Writes are served by the recipient leader, and reads by the donor and recipient leaders.
Once \texttt{TXN\_START($G_i$)} is committed (step~\mstep{2}), writes to its slots are redirected to the recipient shard group, processed by the recipient leader, and committed through its replicated log.
The donor applies no further writes to $G_i$, so no write-back or reconciliation is needed after the migration, and stale-entry detection and tombstones keep the migrated data from overwriting these writes.
As in Fulva~\cite{fulva}, writes go to the recipient and reads use double reads during the migration.
Until \texttt{TRANSFER\_DURABLE($G_i$)} commits, clients use double reads, so that a read never depends on data the recipient may not yet hold.
Once it commits, the recipient shard group serves $G_i$ on its own. By the durability invariant, any newly elected leader can obtain the data from its shard group.

\subsection{Failure handling}
\label{sec:failure-handling}
\label{sec:correctness}

A crash erases everything a node holds only in memory, such as landing buffers and the data in them, chains, and in-flight RDMA writes.
\OurImplementation{} therefore recovers from two kinds of state.
\emph{Durable state} is the set of committed log entries of Table~\ref{tab:log-entries}.
It tells a newly elected leader which groups are in flight, namely those with a start entry but no matching done entry.
\emph{Soft state} records which replicas currently hold the data of a group, and the status of the chain that transfers it.
It is not logged, but a leader can rebuild it from the live replicas.
Recovery always moves forward. A new leader resumes each group in flight from its last committed entry and never rolls it back.
Sending the data of a group again is harmless, because the index update discards entries that are not newer than what the recipient holds.

\paragraph{Soft state.}
Each recipient replica $R_j$ keeps a local flag $\texttt{TRANSFERRED}[R_j][G_i]$, which says that $R_j$ holds the whole of $G_i$.
One-sided RDMA writes do not notify the replica they write to, so the flag is set by the message that follows them.
The leader sets its flag when it receives \texttt{RDMA\_DONE($G_i$)} from the donor, and a follower sets its flag when it receives \texttt{CHAIN\_DONE($G_i$)} from its predecessor in the chain (Section~\ref{sec:chain-replication}).
A leader that fetches $G_i$ from a replica sets its flag once the fetch completes.
Each replica also records the status of the link to its successor in the chain, that is, whether the timeout for its successor has expired or its RDMA connection to it has broken (Section~\ref{sec:chain-replication}).
Each replica reports its flag and the status of its link in every reply to the leader's \texttt{AppendEntries} messages, so the soft state needs no extra messages.
The leader uses the link status to repair the chain, and it reads the flags in two places. It commits \texttt{TRANSFER\_DURABLE($G_i$)} once a majority has set its flag, and it answers \texttt{QUERY\_STATE($G_i$)} from its own flag and those of the other replicas.
A newly elected recipient leader therefore knows from its own flag whether it holds the data, and it learns the flags of the other replicas within one round of \texttt{AppendEntries}.
A replica that crashes restarts with the flag cleared, since its buffers were in memory, so a flag never claims data that no longer exists.

\begin{algorithm}[t]
\small
\caption{Donor Leader Recovery}
\label{alg:donor-recovery}
\begin{algorithmic}[1]

\Upon{\msg{leader\_election\_complete}}
  \If{\textbf{not} (self.group\_type is DONOR \textbf{and} self is LEADER)}
    \State \textbf{return}
  \EndIf

  \If{$\neg$ \fn{Has\_Active\_Migration\_Txn}()}
    \State \textbf{return}
  \EndIf

  \phasesep{Load migration progress from the log}

  \State started\_groups $\leftarrow$ \fn{Log\_Query\_Gi\_With\_State}(\msg{started})
  \State done\_groups $\leftarrow$ \fn{Log\_Query\_Gi\_With\_State}(\msg{done})

  \For{\textbf{each} $G_i$ \textbf{in} started\_groups}

    \If{$G_i \in$ done\_groups}
      \State \fn{Mark\_Group\_Complete}($G_i$) \Comment{completed before failure}
      \State \textbf{continue}
    \EndIf

    \phasesep{Query recipient for $G_i$ state}

    \State \fn{Send\_RPC}(recipient\_leader, \msg{query\_state}($G_i$))
    \State state $\leftarrow$ \fn{Wait\_Reply}(recipient\_leader)

    \Statex

    \If{state $=$ \msg{durable}}
      \State \fn{Log\_Append}(\msg{done}($G_i$)) \Comment{finalize locally}
      \State \fn{Log\_Commit}()

    \ElsIf{state $=$ \msg{transferred}}
      \State \fn{Wait\_For\_Durable}($G_i$) \Comment{timeout: re-query}

    \ElsIf{state $=$ \msg{missing}}
      \State \fn{Send\_RPC}(recipient\_leader, \msg{prepare}($G_i$))
      \State \fn{Start\_RDMA\_Transfer}($G_i$) \Comment{resend}
      \State \fn{Wait\_For\_Durable}($G_i$) \Comment{timeout: re-query}

    \Else
      \State \fn{Defer\_Gi}($G_i$) \Comment{timeout: re-query later}
    \EndIf

  \EndFor

  \phasesep{Global transaction completion}

  \If{\fn{All\_Groups\_Done}()}
    \State \fn{Log\_Append}(\msg{txn\_migration\_done})
    \State \fn{Log\_Commit}()
  \EndIf

\EndUpon

\end{algorithmic}
\end{algorithm}

\begin{algorithm}[t]
\small
\caption{Recipient Leader Recovery}
\label{alg:recipient-leader-recovery}
\begin{algorithmic}[1]

\Upon{\msg{leader\_election\_complete}}
  \If{self.group\_type $\neq$ \texttt{RECIPIENT} \textbf{or} self $\neq$ \texttt{LEADER}}
    \State \textbf{return}
  \EndIf

  \phasesep{Notify the donor of completed groups}

  \For{\textbf{each} $G_i$ with \msg{index\_update} committed}
    \State \fn{Send\_RPC}(donor\_leader, \msg{g\_i\_done}($G_i$)) \Comment{may repeat}
  \EndFor

  \phasesep{Resume incomplete groups}

  \State active\_groups $\leftarrow$ groups with \msg{txn\_start}, no \msg{index\_update}
  \State \textbf{wait} one heartbeat for replica reports

  \For{\textbf{each} $G_i$ \textbf{in} active\_groups}

    \phasesep{Obtain the data first}

    \If{$\texttt{TRANSFERRED}[\texttt{self}][G_i] = \texttt{false}$} \Comment{own flag}
      \If{$\exists R_j$: $\texttt{TRANSFERRED}[R_j][G_i] = \texttt{true}$}
        \State \fn{Fetch\_From\_Replica}($R_j$, $G_i$)
      \Else
        \State \fn{Send\_RPC}(donor\_leader, \msg{resend}($G_i$))
        \State \textbf{wait until} \msg{rdma\_done} of $G_i$ \Comment{timeout: re-request}
      \EndIf
      \State $\texttt{TRANSFERRED}[\texttt{self}][G_i] \leftarrow \texttt{true}$
    \EndIf

    \phasesep{Resume the intra-recipient phase}

    \State chain $\leftarrow$ \fn{Build\_Chain}(live replicas lacking $G_i$)
    \State \fn{Async\_Replicate\_To}(chain, $G_i$)
    \State \fn{Execute\_Update\_Index}($G_i$) \Comment{in parallel}
    \State \textbf{wait until} majority of $\texttt{TRANSFERRED}[R_i][G_i] = \texttt{true}$
    \Statex \hspace{\algorithmicindent}\hspace{\algorithmicindent} \textbf{and} local index update of $G_i$ is done
    \State \fn{Log\_Append}(\msg{index\_update}($G_i$)); \fn{Log\_Commit}()
    \State \fn{Send\_RPC}(donor\_leader, \msg{g\_i\_done}($G_i$))

  \EndFor
\EndUpon

\Statex

\Upon{\fn{Receive} \msg{query\_state}($G_i$) \textbf{from} donor\_leader}
  \If{\msg{index\_update}($G_i$) is committed}
    \State \fn{Reply}(\msg{gi\_done})
  \ElsIf{$\exists$ live $R_j$: $\texttt{TRANSFERRED}[R_j][G_i] = \texttt{true}$}
    \State \fn{Reply}(\msg{transferred}) \Comment{may go stale}
  \Else
    \State \fn{Reply}(\msg{missing})
  \EndIf
\EndUpon

\end{algorithmic}
\end{algorithm}

\iffalse
\begin{figure}[t]
  \centering
  \resizebox{0.9\columnwidth}{!}{\begin{tikzpicture}[x={\columnwidth/252},y=-1pt,
    font=\fontsize{7.5}{9}\fontfamily{phv}\selectfont,
    >={Stealth[length=3.5pt,width=2.8pt]},
    sec/.style={anchor=base west,inner sep=0pt,
                font=\fontsize{7.5}{9}\fontfamily{phv}\bfseries},
    rowlbl/.style={anchor=east,inner sep=0pt,
                   font=\fontsize{7}{8}\fontfamily{phv}\selectfont},
    entry/.style={draw=ink,line width=0.5pt,fill=entfill,rounded corners=1pt,
                  minimum width=30pt,minimum height=8pt,inner sep=0pt},
    noentry/.style={entry,fill=white,draw=lifegray,dash pattern=on 1pt off 1pt},
    flag/.style={draw=ink,line width=0.5pt,fill=ink,minimum size=7pt,inner sep=0pt},
    noflag/.style={flag,fill=white},
    reply/.style={draw=ink,line width=0.6pt,fill=white,rounded corners=1.5pt,
                  minimum width=52pt,minimum height=13pt,inner sep=1pt,
                  font=\fontsize{7}{8}\fontfamily{phv}\bfseries},
    flow/.style={->,draw=ink,line width=0.6pt},
    note/.style={inner sep=0.5pt,font=\fontsize{6}{7}\fontfamily{phv}\itshape},
  ]
  \definecolor{ink}{RGB}{34,34,34}
  \definecolor{lifegray}{RGB}{165,165,165}
  \definecolor{entfill}{gray}{0.80}
  \definecolor{band}{gray}{0.95}
  \def\Gi{\textit{G}\textsubscript{\textit{i}}}
  \def\us{{\fontencoding{T1}\selectfont\_}}%

  \def\xl{88}
  \def\xa{113} \def\xb{171} \def\xc{229}
  \def\ylog{9}  \def\yts{22} \def\ytd{35}
  \def\ysoft{53} \def\ya{66} \def\yb{79} \def\yc{92}
  \def\yrep{118}

  \fill[band,rounded corners=2pt] (0,0) rectangle (252,42);
  \fill[band,rounded corners=2pt] (0,45) rectangle (252,99);

  \node[sec] at (2,\ylog) {Recipient Raft log};
  \node[rowlbl] at (\xl,\yts) {TXN\us{}START(\Gi)};
  \node[rowlbl] at (\xl,\ytd) {TRANSFER\us{}DURABLE(\Gi)};
  \node[entry]   at (\xa,\yts) {}; \node[entry]   at (\xa,\ytd) {};
  \node[entry]   at (\xb,\yts) {}; \node[noentry] at (\xb,\ytd) {};
  \node[entry]   at (\xc,\yts) {}; \node[noentry] at (\xc,\ytd) {};

  \node[sec] at (2,\ysoft) {Soft state (flag)};
  \node[rowlbl] at (\xl,\ya) {R\textsubscript{1} (leader)};
  \node[rowlbl] at (\xl,\yb) {R\textsubscript{2}};
  \node[rowlbl] at (\xl,\yc) {R\textsubscript{3}};
  \node[flag] at (\xa,\ya) {}; \node[flag] at (\xa,\yb) {}; \node[noflag] at (\xa,\yc) {};
  \node[noflag,dash pattern=on 1pt off 1pt] (l) at (\xb,\ya) {};
  \node[flag] (r) at (\xb,\yb) {}; \node[noflag] at (\xb,\yc) {};
  \draw[flow] (r.east) to[out=0,in=0,looseness=2.2]
    node[note,right] {fetch} (l.east);
  \node[noflag] at (\xc,\ya) {}; \node[noflag] at (\xc,\yb) {}; \node[noflag] at (\xc,\yc) {};

  \node[rowlbl] at (80,\yrep) {Reply};
  \node[reply] (d) at (\xa,\yrep) {DURABLE};
  \node[reply] (t) at (\xb,\yrep) {TRANSFERRED};
  \node[reply] (m) at (\xc,\yrep) {MISSING};
  \foreach \x/\n in {\xa/d,\xb/t,\xc/m} \draw[flow] (\x,101) -- (\n.north);
\end{tikzpicture}
}
  \caption{Reply of the recipient leader to \texttt{QUERY\_STATE($G_i$)}.
  Gray boxes are committed log entries, and filled squares are set \texttt{TRANSFERRED} flags.}
  \label{fig:query_state}
\end{figure}
\fi
\paragraph{Donor recovery.}
A newly elected donor leader resumes each group that has \texttt{START($G_i$)} but no \texttt{DONE($G_i$)} (Algorithm~\ref{alg:donor-recovery}).
It asks the recipient leader for the state of the group with a \texttt{QUERY\_STATE($G_i$)} RPC, and the recipient answers from its soft state and its log.
The reply is \texttt{DURABLE} if \texttt{TRANSFER\_DURABLE($G_i$)} is committed.
Otherwise, the reply is \texttt{TRANSFERRED} if the recipient leader holds the data, that is, if its own flag $\texttt{TRANSFERRED}[R_j][G_i]$ is set.
If it does not, it asks the other replicas, fetches the data from one whose flag is set, and then replies \texttt{TRANSFERRED}.
If no live replica holds the data, the reply is \texttt{MISSING}.
The donor commits \texttt{DONE($G_i$)} on \texttt{DURABLE}.
On \texttt{MISSING}, it first transfers the group again, starting from \texttt{prepare}.
On any reply other than \texttt{DURABLE}, it keeps sending \texttt{QUERY\_STATE($G_i$)} until it receives \texttt{DURABLE}.
If the replicas holding the data crash after a \texttt{TRANSFERRED} reply, the recipient leader asks for the data again, and the donor can always serve it, because it keeps its copy of $G_i$ until \texttt{DONE($G_i$)} commits.
Every new transfer uses newly allocated landing buffers with new RKeys, so late RDMA writes from an earlier attempt fail instead of overwriting them.

\paragraph{Recipient recovery.}
A newly elected recipient leader resumes each group that has \texttt{TXN\_START($G_i$)} but no \texttt{TRANSFER\_DURABLE($G_i$)} (Algorithm~\ref{alg:recipient-leader-recovery}).
A group whose \texttt{TRANSFER\_DURABLE($G_i$)} is committed is complete on the recipient, but the old leader may have crashed before telling the donor.
The new leader does nothing for such a group, because the donor learns the outcome from the reply to its \texttt{QUERY\_STATE($G_i$)}.
For a group in flight, the leader must first hold the data itself.
If its transferred flag is not set, it fetches the data from a replica whose flag is set.
If no live replica holds the data, it asks the donor to send it again, which is the usual case when the crash interrupted a transfer, since the followers had not yet received the whole group.
The leader then resumes the intra-recipient phase. It replicates the data to the replicas that lack it, updates its index, and commits \texttt{TRANSFER\_DURABLE($G_i$)} once a majority holds the data.
It obtains the data before it waits for a majority, because if every replica that held the data has crashed, no majority can form until the data arrives again.

\paragraph{Timeouts.}
Failures are detected through timeouts or broken connections (Section~\ref{sec:fault-model}), so every wait that depends on another node has a timeout.
A donor leader that gets no \texttt{DURABLE} after sending \texttt{RDMA\_DONE($G_i$)} sends \texttt{QUERY\_STATE($G_i$)} to the recipient leader, finding the new leader if it has changed.
A donor leader that gets no reply to \texttt{QUERY\_STATE($G_i$)}, or no \texttt{DURABLE} after a \texttt{TRANSFERRED} reply or a new transfer, queries the recipient leader again.
A recipient leader that asked for the data and gets no \texttt{RDMA\_DONE} asks again.
A newly elected recipient leader waits one heartbeat interval for the flags of the replicas and treats a replica that does not answer as lacking the data.
Every action taken on a timeout is a query or a new transfer, both of which can be repeated safely, and every decision comes from committed log entries.
A timeout that fires too early therefore costs at most an extra transfer, and one that fires too late only delays recovery.

\iffalse
Fig.~\ref{fig:recipient_recovery} shows a recipient leader crash: B1 crashes after receiving $G_i$ but before replicating it, so the new leader B2 finds $G_i$ in flight without holding its data, obtains it again from the donor, and completes the migration.

\begin{figure}[t]
  \centering
  % Fig. 2 -- recovery of one group after the recipient leader crashes.
% Same conventions as figs/1.tex: 252-unit-wide grid (x unit = \columnwidth/252),
% y in pt growing downwards, labels at 7.5pt. \input it directly.
\begin{tikzpicture}[x={\columnwidth/252},y=-1pt,
    font=\fontsize{7.5}{9}\fontfamily{phv}\selectfont,
    >={Stealth[length=4pt,width=3.2pt]},
    lifeline/.style={draw=lifegray,line width=0.8pt},
    deadline/.style={draw=lifegray,line width=0.8pt,dash pattern=on 1pt off 2pt},
    shard/.style={draw=ink,line width=0.6pt,fill=white,rounded corners=1.5pt,
                  minimum width=24pt,minimum height=12pt,inner sep=0pt},
    leader/.style={shard,fill=ink,text=white,font=\fontsize{7.5}{9}\fontfamily{phv}\bfseries},
    msg/.style={->,draw=ink,line width=0.6pt},
    bulk/.style={->,draw=ink,line width=1.8pt,>={Stealth[length=5pt,width=5pt]}},
    notify/.style={->,draw=ink,line width=0.6pt,dash pattern=on 2.2pt off 1.4pt},
    crash/.style={draw=ink,line width=1.6pt,line cap=round},
    lbl/.style={inner sep=1pt,text=ink},
    phase/.style={rounded corners=2pt},
    phaselbl/.style={anchor=north west,inner sep=1.5pt,outer sep=0pt,
                     font=\fontsize{7.5}{9}\fontfamily{phv}\bfseries},
  ]
  \definecolor{ink}{RGB}{34,34,34}
  \definecolor{lifegray}{RGB}{165,165,165}
  \definecolor{llfill}{gray}{0.93}  % before the crash
  \definecolor{irfill}{gray}{0.86}  % recovery
  \def\Gi{\textit{G}\textsubscript{\textit{i}}}
  \def\us{{\fontencoding{T1}\selectfont\_}}% thin underscore glyph
  % lettered step badge, same geometry as the numbered badges of Fig. 1
  \def\rstep#1{\tikz[baseline={([yshift=-2.7pt]s.center)}]\node[circle,fill=ink,text=white,inner sep=0pt,minimum size=9.6pt,text height=5.1pt,text depth=0.7pt,font=\fontsize{7.2}{7.2}\fontfamily{phv}\bfseries\selectfont](s){#1};}
  \def\xmark{\tikz[baseline=-2.7pt]{\draw[crash] (-3pt,-3pt) -- (3pt,3pt); \draw[crash] (-3pt,3pt) -- (3pt,-3pt);}}

  % lifeline x positions
  \def\Ah{26}                              % donor leader A3
  \def\Bo{100} \def\Bt{163} \def\Bh{226}  % recipient: B1 (old leader), B2, B3
  \def\crashy{97}
  \def\bottom{302}

  % phase backgrounds
  \fill[phase,llfill] (0,50) rectangle (252,106);
  \fill[phase,irfill] (0,111) rectangle (252,282);

  % lifelines; B1 is dotted after its crash
  \foreach \x in {\Ah,\Bt,\Bh} \draw[lifeline] (\x,33) -- (\x,\bottom);
  \draw[lifeline] (\Bo,33) -- (\Bo,\crashy);
  \draw[deadline] (\Bo,\crashy) -- (\Bo,\bottom);

  % replica groups and shards
  \node[font=\fontsize{7.5}{9}\fontfamily{phv}\bfseries,anchor=base] at (\Ah,6) {Donor};
  \node[font=\fontsize{7.5}{9}\fontfamily{phv}\bfseries,anchor=base] at (\Bt,6) {Recipient shard group};
  \draw[ink,line width=0.5pt] (\Ah-14,15) -- ++(0,-3) -- (\Ah+14,12) -- ++(0,3);
  \draw[ink,line width=0.5pt] (\Bo-14,15) -- ++(0,-3) -- (\Bh+12,12) -- ++(0,3);
  \node[leader] at (\Ah,27) {A3};
  \node[leader] at (\Bo,27) {B1};
  \node[shard]  at (\Bt,27) {B2};
  \node[shard]  at (\Bh,27) {B3};
  \node[lbl,fill=white,font=\fontsize{7}{8}\fontfamily{phv}\itshape] at (\Ah,41) {leader};
  \node[lbl,fill=white,font=\fontsize{7}{8}\fontfamily{phv}\itshape] at (\Bo,41) {leader};

  % before the crash: B1 receives G_i but crashes before replicating it
  \node[phaselbl,fill=llfill] at (2,52) {Before the crash};
  \draw[bulk] (\Ah,76) -- node[lbl,above=2pt,fill=llfill] {transfer} (\Bo,76);
  \draw[msg]  (\Ah,89) -- node[lbl,above=2pt,fill=llfill] {done}     (\Bo,89);
  \draw[crash] (\Bo-4.2,\crashy-4) -- (\Bo+4.2,\crashy+4);
  \draw[crash] (\Bo-4.2,\crashy+4) -- (\Bo+4.2,\crashy-4);
  \node[lbl,right=5pt,fill=llfill,font=\fontsize{7.5}{9}\fontfamily{phv}\bfseries] at (\Bo,\crashy) {B1 crashes};

  % recovery, driven by the new recipient leader B2
  \node[phaselbl,fill=irfill] at (2,113) {Recovery};
  \node[leader,minimum width=30pt] at (\Bt,126) {leader};
  \node[lbl,right=3pt,fill=irfill,anchor=west] at (\Bt+15,126) {\rstep{a}\hspace{2pt}elected};

  \draw[msg,rounded corners=2pt] (\Bt,141) -- ++(12,0) -- ++(0,9) -- (\Bt,150);
  \node[lbl,right=3pt,fill=irfill,align=left,anchor=west] at (\Bt+12,145.5)
    {\rstep{b}\hspace{2pt}\Gi{} in flight,\\\hspace{11.6pt}data missing};

  \draw[msg]  (\Bt,170) -- node[lbl,above=2pt,fill=irfill] {\rstep{c}\hspace{2pt}resend \Gi} (\Ah,170);
  \draw[msg]  (\Ah,183) -- node[lbl,above=2pt,fill=irfill] {prepare}  (\Bt,183);
  \draw[msg]  (\Bt,196) -- node[lbl,above=2pt,fill=irfill] {MRs}      (\Ah,196);
  \draw[bulk] (\Ah,211) -- node[lbl,above=2pt,fill=irfill] {\rstep{d}\hspace{2pt}transfer} (\Bt,211);
  \draw[msg]  (\Ah,224) -- node[lbl,above=2pt,fill=irfill] {done}     (\Bt,224);

  \draw[bulk] (\Bt,241) -- node[lbl,above=2pt,align=center,font=\fontsize{7.5}{8}\fontfamily{phv}\selectfont] {\rstep{e}\hspace{2pt}chain\\replication} (\Bh,241);
  \draw[msg]  (\Bh,254) -- node[lbl,above=2pt,fill=irfill] {ack (2/3)} (\Bt,254);
  \node[lbl,above=2pt,fill=irfill] at ({(\Bt+\Bh)/2},274) {\rstep{f}\hspace{2pt}INDEX\us UPDATE};
  \draw[notify] (\Bt,274) -- (\Bh,274);

  % B2 reports completion; the donor commits DONE(G_i)
  \node[lbl,above=2pt,fill=white] at ({(\Ah+\Bt)/2},297) {\rstep{g}\hspace{2pt}DONE(\Gi)};
  \draw[msg] (\Bt,297) -- (\Ah,297);

  % legend: one centred row, equal gaps
  \def\key#1{\tikz[baseline=-2.7pt]\draw[#1] (0pt,0pt) -- (16pt,0pt);\hspace{3pt}}
  \node[inner sep=0pt,anchor=base] at (126,318) {%
    \key{msg}RPC\hspace{11pt}%
    \key{notify}Log entry\hspace{11pt}%
    \key{bulk}RDMA write\hspace{11pt}%
    \xmark\hspace{3pt}Crash};
\end{tikzpicture}%

  \caption{Recovery of group $G_i$ after the recipient leader crashes.
  (a)~B2 is elected leader.
  (b)~Its log holds \texttt{TXN\_START($G_i$)} but not \texttt{TRANSFER\_DURABLE($G_i$)}, and B2 lacks the data.
  (c--d)~The donor resends $G_i$ to B2.
  (e)~B2 replicates it to B3, reaching a majority.
  (f)~B2 commits \texttt{TRANSFER\_DURABLE($G_i$)}.
  (g)~The donor commits \texttt{DONE($G_i$)}.}
  \label{fig:recipient_recovery}
\end{figure}
\fi

\subsection{Chain replication}
\label{sec:chain-replication}

The recipient leader replicates the data of $G_i$ to its followers with chain replication over RDMA, in parallel with its own index update of $G_i$ (Section~\ref{sec:index_updating}).
The leader orders the live replicas into a chain and writes the data to the first one.
Each replica that has received the data forwards it to its successor and signals it with a \texttt{CHAIN\_DONE} message~(step~\mstep{4}).
The chain is set up before the migration starts, off its critical path.
The recipient leader opens an RDMA connection to every follower and registers landing buffers on each of them, and the followers connect to each other, so any replica can forward the data to any other.
Each transaction gets its own transaction ID and its own landing buffers on every replica, so the groups of concurrent donors never overwrite each other.

\paragraph{Transfer state.}
The leader learns the state of the transfer without extra messages.
Every replica piggybacks it on its replies to the leader's \texttt{AppendEntries} messages, reporting which landing buffers of $G_i$ it holds and whether its transferred flag is set.
A replica sets the flag only after the whole of $G_i$ has landed, that is, after it receives \texttt{RDMA\_DONE} from the donor or \texttt{CHAIN\_DONE} from its predecessor, because one-sided RDMA writes do not notify it.
The leader therefore knows at any time how far the data has travelled down the chain.
It learns that a majority holds $G_i$ once $f$ followers have reported it, without waiting for the rest of the chain.
Replication to this majority is synchronous. The remaining replicas receive the data asynchronously from the ongoing chain.

\paragraph{Chain repair.}
Each replica in the chain keeps a timeout for its successor.
The timeout starts when the replica forwards $G_i$ to its successor, and stops when the successor acknowledges that it holds $G_i$.
If the timeout expires, the replica reports it to the leader in its next reply to the leader's \texttt{AppendEntries} messages.
The leader then moves the successor to the end of the chain.
It tells the closest replica before it that holds $G_i$ to forward the data to the next replica in the new order.
This is the leader itself if no follower before the moved replica holds $G_i$.
The data does not change during the transfer, so the forwarding replica sends the whole of $G_i$ again and needs no record of what it sent before.
A replica whose RDMA connection to its successor breaks reports it in the same way, without waiting for the timeout, and the leader moves that successor to the end of the chain as well.
A replica that crashed and rejoins reports a cleared flag, and the leader adds it at the end of the chain.

Chain repair affects only liveness.
The leader counts a replica as holding $G_i$ only when that replica reports it, so a timeout never makes \texttt{TRANSFER\_DURABLE($G_i$)} commit early.
A timeout that fires too early only moves a slow replica to the end of the chain, where it still receives the data.
If no majority can be reached, \texttt{TRANSFER\_DURABLE($G_i$)} never commits, and the migration fails visibly instead of silently losing data.

